\section{位置と道のり}\label{節:運動と積分}

\sectionref{節:符号と道のり}では,速度が一定の区間をつないだ運動で,速度に時間を掛けて足すと位置の変化,速さに時間を掛けて足すと道のりになることを見た.では,速度が連続的に変わる場合はどうなるだろうか.

一直線上の運動で,時刻 $t$ の位置を $s(t)$,速度を $v(t)=s'(t)$,加速度を $a(t)=v'(t)$ とし,\,$v$ は連続であるとする.時刻 $t_0$ から $t_1$ までを細かく分けると,各小区間では速度がほぼ一定とみなせる.長さ $\Delta t$ の小区間で,位置の変化はほぼ $v(t)\Delta t$,進んだ距離はほぼ $|v(t)|\Delta t$ である.これらの和の極限として,位置の変化も道のりも定積分に帰着する.

位置の変化については,微積分の基本定理から直接
\[
s(t_1)-s(t_0)=\int_{t_0}^{t_1}v(t)\dd t
\]
を得る.左辺は位置の差,すなわち\textbf{変位}である.出発点に戻れば変位は $0$ になるが,途中で移動した距離が $0$ という意味ではない.

速度が正なら正の向きに,負なら負の向きに進む.向きによらず移動した距離を足すには,速さ $|v(t)|$ を積分して
\begin{equation}\label{式:道のり}
\text{道のり}=\int_{t_0}^{t_1}|v(t)|\dd t
\end{equation}
とする.速度には向きに対応する符号があり,速さは負にならない.この違いは,符号付きの定積分と面積の違いと同じである.

速度が $0$ になる時刻は,進む向きが変わる候補である.ただし,\,$0$ になるだけでは向きが変わるとは限らない.たとえば $v(t)=(t-1)^2$ は $t=1$ で $0$ になるが,前後で負にならない.道のりを計算するときは,速度が $0$ になる時刻で区間を分けたうえで,各区間の符号を確かめる.

\begin{bookpoint}
\textbf{POINT}\par
速度が $0$ になる時刻で区間を分け,各区間で $\int v\dd t$(位置の差)の絶対値を求めて足す.
\end{bookpoint}

\begin{example}{変位と道のり}
時刻 $0\le t\le4$ において,速度が $v(t)=(t-1)(t-3)$ で与えられる.初期位置は $s(0)=0$ とする.\,$t=4$ の位置,道のり,平均速度,平均の速さを求めよ.
\end{example}

\paragraphbf{方針}
速度を積分して位置 $s(t)$ を求める.\,$v(t)=(t-1)(t-3)$ は $1<t<3$ で負,それ以外では $0$ 以上なので,道のりは $[0,1]$,\,$[1,3]$,\,$[3,4]$ に分けて,各区間の位置の差の絶対値を足す.

\paragraphbf{解答}
位置 $\frac43$,道のり $4$,平均速度 $\frac13$,平均の速さ $1$

\paragraphbf{解法}
$v(t)=t^2-4t+3$ を積分し,\,$s(0)=0$ を使うと
\[
s(t)=\frac{t^3}{3}-2t^2+3t
\]
である.よって
\begin{align*}
s(4)&=\frac{64}{3}-32+12\\
 &=\frac43.
\end{align*}
次に速度の符号を調べる.\,$v$ は $0\le t<1$ で正,\,$1<t<3$ で負,\,$3<t\le4$ で正である.向きを変える時刻の位置は
\[
s(1)=\frac43,\qquad s(3)=0
\]
なので,道のりは
\begin{align*}
&\left|s(1)-s(0)\right|+\left|s(3)-s(1)\right|+\left|s(4)-s(3)\right|\\
&=\frac43+\frac43+\frac43\\
&=4
\end{align*}
となる.これは $\int_0^1v\dd t-\int_1^3v\dd t+\int_3^4v\dd t$ を計算したことと同じである.全時間は $4$ だから,平均速度と平均の速さはそれぞれ
\[
\frac{s(4)-s(0)}{4}=\frac13,\qquad
\frac{4}{4}=1
\]
となる.

\paragraphbf{速さが増える区間}
例題の加速度は $a(t)=2t-4$ である.加速度が正だからといって,速さが増えるとは限らない.\,$v(t)\ne0$ の時刻では
\[
\frac{\dd}{\dd t}|v(t)|=
\begin{cases}
a(t),&v(t)>0,\\
-a(t),&v(t)<0
\end{cases}
\]
であるから,速度と加速度が同符号なら速さが増え,異符号なら減る.例題では,\,$2<t<3$ で加速度は正だが速度は負なので,速さは減っている.逆に $1<t<2$ では速度も加速度も負なので,負の向きに進みながら速くなっているのだ.速さが増えるのは $1<t<2$ と $3<t<4$ である.

\paragraphbf{平面上の運動}
平面上を動く点 $\bigl(x(t),y(t)\bigr)$ では,速度は2つの成分 $\bigl(x'(t),y'(t)\bigr)$ を持ち,速さは三平方の定理から $\sqrt{x'(t)^2+y'(t)^2}$ となる.この速さを時刻 $\alpha$ から $\beta$ まで積分した道のりは,式\eqref{式:媒介曲線の弧長}と同じ式である.つまり,媒介変数 $t$ を時刻とみれば,曲線の長さの公式は動く点の道のりを求める式になっている.ただし運動では,同じ場所を再び通った分も道のりに含める.曲線を重複してたどった長さは,図形の長さとしては数えすぎだが,道のりとしては数えるべき量なのである.
